\bookchapter{Глава 52. Четыре компонента и жизненный цикл Activity}

\chapterintro{Если Compose — «лицо» приложения, то эта глава о его «скелете». В Android есть четыре базовых компонента, из которых собрано любое приложение; а у важнейшего из них, Activity, есть «карта жизненного цикла», которую каждый Android-разработчик обязан вырезать на памяти. Осилив главу, вы наконец поймёте вечные загадки вроде «почему после сворачивания моё приложение забыло всё своё состояние».}

\booksection{Четыре компонента: несущие опоры Android}

Официально в Android определены четыре базовых компонента; каждый отвечает за свой класс задач:

\begin{center}
\begin{tabularx}{\linewidth}{|X|X|X|}
\hline
\textbf{Компонент} & \textbf{Обязанность} & \textbf{Для наглядности} \tabularnewline
\hline
Activity & Видимый экран & Витрина магазина — её видно, с ней взаимодействуют \tabularnewline
Service & «Невидимка», работающая в фоне & Складской рабочий — играет музыку, качает файлы, экран ему не нужен \tabularnewline
BroadcastReceiver & Принимает системные оповещения & Проходная — «рупор» вестей вроде «батарея садится» или «сеть пропала» \tabularnewline
ContentProvider & Делится данными между приложениями & Общий архив — например, открывает контакты другим приложениям \tabularnewline
\hline
\end{tabularx}
\end{center}

\begin{center}\small\color{gray}Четыре компонента Android: краткий обзор\end{center}

В роли новичка 95\% времени вы будете возиться с Activity; Service и BroadcastReceiver понадобятся для музыки и push-уведомлений; ContentProvider нужен реже всего (типовые данные система уже упаковала). И учтите: все четыре компонента обязаны «прописаться» в AndroidManifest.xml — непрописанных система не найдёт.

\booksection{Жизненный цикл Activity: первая «карта сокровищ» Android-разработчика}

Приложение на телефоне могут прервать в любой момент: звонок, пользователь ушёл ответить в WeChat, системе не хватило памяти, и она вас убила… Жизнь Activity — вечное странствие по этой карте:

onCreate (рождение: инициализация)

onStart (виден)

onResume (на переднем плане, взаимодействие)

onPause (частично перекрыт)

onStop (полностью скрыт)

onDestroy (уничтожен)

onRestart / onStart / onResume: при возврате на передний план — обратный путь

\begin{center}\small\color{gray}Рисунок 4. Жизненный цикл Activity (по ходу: верхний ряд — рождение нижний ряд — уход; обратно — тем же путём)\end{center}

\begin{center}
\begin{tabularx}{\linewidth}{|X|X|X|}
\hline
\textbf{Колбэк} & \textbf{Когда срабатывает} & \textbf{Типичное применение} \tabularnewline
\hline
onCreate & При создании Activity (один раз) & Загрузка интерфейса, инициализация данных \tabularnewline
onStart & Интерфейс стал видимым & Запуск нетяжёлой логики \tabularnewline
onResume & Получен фокус, можно взаимодействовать & Запуск анимаций, продолжение воспроизведения \tabularnewline
onPause & Экран частично перекрыт / момент ухода с него & Пауза анимаций, сохранение черновика \tabularnewline
onStop & Полностью невидим & Освобождение тяжёлых ресурсов, пауза загрузок \tabularnewline
onDestroy & Перед уничтожением & Финальная уборка \tabularnewline
\hline
\end{tabularx}
\end{center}

\begin{center}\small\color{gray}Типичное применение шести колбэков жизненного цикла\end{center}

\textbf{Осторожно: поворот экрана «убивает» Activity}\par

Классическое фиаско новичка: повернулся экран — и весь текст в полях ввода исчез. Причина: по умолчанию поворот экрана уничтожает и заново создаёт Activity (проходит цикл onDestroy onCreate). Решение: сохраняйте состояние через rememberSaveable (замена remember в Compose) — при пересоздании он сам вернёт значения. Пригодится и в главах 6, и в главе 7.

\booksection{Intent: «билет» на поездку из приложения и внутрь него}

Intent («намерение») — «билет» для общения компонентов. Хотите перейти с экрана на экран, открыть браузер или вызвать камеру — сначала «купите билет и сядьте в вагон»:

\begin{lstlisting}[language=Kotlin]
// Способ 1: явный Intent — прямо указываем, на какой экран идти
val intent = Intent(this, DetailActivity::class.java)
intent.putExtra("name", Иван)     // везём с собой «багаж» (данные)
startActivity(intent)

// Способ 2: неявный Intent — говорим, что сделать, система ищет исполнителя
val web = Intent(Intent.ACTION_VIEW, Uri.parse("https://developer.android.com"))
startActivity(web)   // система откроет браузер / или предложит выбрать
\end{lstlisting}

У двух способов разные роли: явный (explicit) — для переходов между экранами своего приложения; неявный (implicit) — чтобы обратиться к возможностям системы или чужих приложений (открыть сайт, позвонить, поделиться). В Compose-проекте возьмите Context через LocalContext.current прямо в Composable — и вызывайте startActivity.

\booksection{Несколько экранов и навигация}

В настоящих приложениях экранов много. Официально рекомендованное для Compose навигационное решение — Navigation Compose (библиотека androidx.navigation):

\begin{lstlisting}[language=Kotlin]
val navController = rememberNavController()

NavHost(navController, startDestination = "home") {
    composable("home") { HomeScreen(navController) }
    composable("detail/{name}") { entry ->
        val name = entry.arguments?.getString("name")
        DetailScreen(name ?: "Неизвестно")
    }
}

// Переход внутри экрана:
navController.navigate("detail/Иван")     // с аргументами
navController.popBackStack()              // назад
\end{lstlisting}

\begin{tcbinfo}{Совет: экранов становится много — что делать?}
Когда проект вырастет, соберите все маршруты экранов в одном файле (например, AppNavGraph.kt) и управляйте ими централизованно — не раскидывайте логику «куда перейти» по разным экранам. Это первый шаг от «работает» к «легко поддерживать».

\end{tcbinfo}

\booksection{Заодно знакомимся: что такое Fragment (и почему книга его не преподаёт)}

В старых учебниках вы непременно наткнётесь на Fragment — «переиспользуемый кусочек интерфейса» внутри Activity, абсолютную звезду старой многоэкранной архитектуры. Но в эпоху Compose для экранов хватает Composable-функций и маршрутов Navigation, и официально роль Fragment ослабла. Поэтому основная линия книги его обходит.

Мой совет: в новых проектах сами Fragment не тащите; но при чтении старого кода или вопросе на собеседовании достаточно знать, что это «кусочек интерфейса, встроенный в Activity». Учебные материалы всегда берите не старше 2023 года — «старый календарь» Android-стека обновляется быстрее, чем кажется.

\begin{tcbpractice}{Практика: научите счётчик «помнить»}
Доработайте счётчик из главы 51: (1) замените remember на rememberSaveable и, повернув экран, проверьте, что число на месте; (2) добавьте кнопку «О программе», которая через Navigation ведёт на второй экран с номером версии; (3) добавьте в onCreate/onPause обоих экранов вывод в журнал (Log.d) и понаблюдайте в Logcat за порядком вызовов жизненного цикла. Сделаете эти три шага — и жизненный цикл перестанет быть абстракцией.

\end{tcbpractice}

\begin{tcbsummary}{Коротко о главном}
\begin{itemize}
  \item Четыре компонента: Activity (витрина), Service (фон), BroadcastReceiver (рупор), ContentProvider (архив)
  \item Шесть колбэков жизненного цикла: onCreate onStart onResume onPause onStop onDestroy
  \item Поворот экрана уничтожает и пересоздаёт Activity: спасайте состояние через rememberSaveable
  \item Intent — «билет»: явный ведёт на свои экраны, неявный зовёт возможности системы
  \item Несколько экранов — через Navigation Compose: NavHost + маршруты composable
\end{itemize}

\end{tcbsummary}
